\documentclass[11pt]{article}% Shared preamble for the rigor documents (Lamport-style structured proofs).  \input this file after \documentclass.
\usepackage[T1]{fontenc}\usepackage{lmodern}\usepackage{microtype}
\usepackage[margin=1in]{geometry}
\usepackage{amsmath,amssymb,amsthm,mathtools}
\usepackage{xcolor}\usepackage{enumitem}\usepackage{booktabs}\usepackage{graphicx}\usepackage{hyperref}
\usepackage[most]{tcolorbox}
\definecolor{navy}{HTML}{17365D}\definecolor{teal}{HTML}{117A8B}\definecolor{warn}{HTML}{A33A2B}\definecolor{okgreen}{HTML}{2E7D4F}
\hypersetup{colorlinks=true,linkcolor=navy,citecolor=teal,urlcolor=teal}
\theoremstyle{plain}
\newtheorem{theorem}{Theorem}[section]\newtheorem{lemma}[theorem]{Lemma}\newtheorem{proposition}[theorem]{Proposition}\newtheorem{corollary}[theorem]{Corollary}
\theoremstyle{definition}\newtheorem{definition}[theorem]{Definition}\newtheorem{remark}[theorem]{Remark}\newtheorem{claim}[theorem]{Claim}\newtheorem{issue}[theorem]{Issue}
% ---------- Lamport structured proofs ----------
% \lstep{level}{number}{assertion}   -> indented "<level>number. assertion"
% \lproof                             -> "PROOF:" line (start of the sub-proof of the preceding step)
% \lby{justification}                 -> "BY ..." leaf justification for the preceding step
% \lqed{level}{number}                -> QED step of a sub-proof
% keywords: \lassume \lprove \lsuffices \lcase \ldefine \lpick \lnew
\newlength{\lindent}\setlength{\lindent}{1.7em}
\newcommand{\lstep}[3]{\par\smallskip\noindent\hspace*{\dimexpr\numexpr#1-1\relax\lindent\relax}{\color{navy}$\langle#1\rangle#2$.}\ #3}
\newcommand{\lproof}{\par\noindent\hspace*{\lindent}{\scshape Proof:}}
\newcommand{\lby}[1]{\par\noindent\hspace*{\lindent}{\scshape By}\ #1.}
\newcommand{\lqed}[2]{\par\smallskip\noindent\hspace*{\dimexpr\numexpr#1-1\relax\lindent\relax}{\color{navy}$\langle#1\rangle#2$.}\ {\scshape Q.E.D.}}
\newcommand{\lassume}{{\scshape Assume}\ }\newcommand{\lprove}{{\scshape Prove}\ }\newcommand{\lsuffices}{{\scshape Suffices}\ }
\newcommand{\lcase}{{\scshape Case}\ }\newcommand{\ldefine}{{\scshape Define}\ }\newcommand{\lpick}{{\scshape Pick}\ }\newcommand{\lnew}{{\scshape New}\ }
% ---------- verdict boxes ----------
\newtcolorbox{verdictok}{colback=okgreen!8,colframe=okgreen,title={\bfseries Verdict: verified},fonttitle=\small}
\newtcolorbox{verdictgap}{colback=orange!10,colframe=orange!80!black,title={\bfseries Verdict: gap / caveat},fonttitle=\small}
\newtcolorbox{verdictbreak}{colback=warn!10,colframe=warn,title={\bfseries Verdict: proof-breaking issue},fonttitle=\small}
\newtcolorbox{citedbox}[1]{colback=navy!5,colframe=navy!60,title={\bfseries Cited result: #1},fonttitle=\small,breakable}
\setlength{\parindent}{0pt}\setlength{\parskip}{4pt}\allowdisplaybreaks
\newcommand{\cA}{\mathcal A}\newcommand{\cF}{\mathcal F}\newcommand{\cM}{\mathfrak M}\newcommand{\cN}{\mathfrak N}

% extra calligraphic shortcuts (\providecommand: no clash if a document defines its own)
\providecommand{\cB}{\mathcal B}\providecommand{\cC}{\mathcal C}\providecommand{\cD}{\mathcal D}\providecommand{\cE}{\mathcal E}\providecommand{\cG}{\mathcal G}\providecommand{\cH}{\mathcal H}\providecommand{\cI}{\mathcal I}\providecommand{\cJ}{\mathcal J}\providecommand{\cK}{\mathcal K}\providecommand{\cL}{\mathcal L}\providecommand{\cO}{\mathcal O}\providecommand{\cP}{\mathcal P}\providecommand{\cQ}{\mathcal Q}\providecommand{\cR}{\mathcal R}\providecommand{\cS}{\mathcal S}\providecommand{\cT}{\mathcal T}\providecommand{\cU}{\mathcal U}\providecommand{\cV}{\mathcal V}\providecommand{\cW}{\mathcal W}\providecommand{\cX}{\mathcal X}\providecommand{\cY}{\mathcal Y}\providecommand{\cZ}{\mathcal Z}
\newcommand{\Fid}{\operatorname{F}}\newcommand{\Tr}{\operatorname{Tr}}\newcommand{\eps}{\varepsilon}\newcommand{\dd}{\,\mathrm d}

\newcommand{\hh}{\mathfrak h}
\newcommand{\Herm}{\operatorname{Herm}}
\newcommand{\SLD}{\mathcal S_Q}
\newcommand{\Aop}{\mathbb A}
\newcommand{\Stwo}{\mathfrak S_2}
\DeclareMathOperator{\ran}{ran}
\DeclareMathOperator{\dist}{dist}
\newcommand{\Erecc}{E_{\mathrm{rec}}}
\title{Referee report: \texttt{rigor/exact\_optimum\_tangent\_problem.tex}\\
\large (track F1, Phase 5; adversarial check of the exact chart of $\cF$, the tangent
problem, the boundary-direction theorem, the KKT sign criterion and the modular frame)}
\author{Referee REF-P5-1 (Opus)}\date{2026-09-15}
\begin{document}\maketitle
\begin{abstract}
Adversarial report on F1's \texttt{exact\_optimum\_tangent\_problem.tex} (19\,pp, 1135
lines, 2026-09-15).  Every lemma, proposition and theorem receives one verdict
(\textsc{stands} / \textsc{stands with repair} / \textsc{fails}), with the precise step, an
independent recomputation and, where the claim is false, an explicit counterexample.  The
mandated recomputations (Lemmas~2.1--2.2, all blocks of Prop.~3.4, Lemma~3.2 at finite $s$,
Thm.~3.11, Thm.~4.4 including a numerical counterexample search, Prop.~4.6, Prop.~6.2) were
carried out from scratch; numerical spot checks used \texttt{numpy} on random matrices.
Headline: the chart, the first-order blocks, the exact feasibility, the derivative formula,
the equivalence (a)$\Leftrightarrow$(c) and the modular-frame reduction all stand; the
boundary-direction theorem (Thm.~3.11(3)) is \emph{false as stated} and takes with it the
last clause of Lemma~4.2 and the completeness of Thm.~4.4(d); the $\theta$-diagnosis of
\S6.3 is superseded by F3.  \S\ref{sec:table} collects the verdicts and \S\ref{sec:repairs}
the required repairs.
\end{abstract}
\tableofcontents

\section{Scope, method, and what was checked against what}\label{sec:scope}

The refereed document is \texttt{rigor/exact\_optimum\_tangent\_problem.tex} (1135 lines,
2026-09-15, track F1).  It was read in full, in pieces, and every invocation of another
document was checked at the cited place: \texttt{sdp\_dual\_certificate.tex} (``S10'':
(C4), (C5), (N4), (N5), Def.~3.1, Thm.~3.2$\langle1\rangle5$--$\langle1\rangle8$, Cor.~3.3,
Prop.~5.2, Cor.~5.3, Lemma~6.1, Thm.~6.2, \S7, Lemma~11.1, Cor.~11.2),
\texttt{optimality\_second\_order.tex} (``S8'': Thm.~3.1, eq.~(3.9)),
\texttt{optimality\_all\_channels.tex} (``S9'': Thm.~3.5, Thm.~4.1, \S5.2) and
\texttt{exact\_fidelity\_upper\_half.tex} (``E1'': Cor.~7.5).  Numerical spot checks were
run in \texttt{numpy} (random Hermitian/positive matrices, $n\le12$); they are reported
inline and are not part of any proof.  Two external inputs postdating the document are used
where they bear on it: track F2's KKT numerics (\texttt{numerics/optimality\_all/F2\_RESULTS.md})
and track F3's $\theta$ reconciliation (\texttt{F3\_RESULTS.md}, \texttt{f3\_t0.py}).

Adversarial stance: the burden is on the document.  A claim is \textsc{stands} only if the
referee could reproduce it independently; \textsc{stands with repair} if the statement
survives but a step, a hypothesis or a domain must be changed; \textsc{fails} if a
counterexample to the claim \emph{as written} exists.

\section{Conventions: the SLD normalisation and the factor $2$ in S10~(C4)}\label{sec:norm}

\begin{lemma}[Recomputation of Lemma~1.1]\label{lem:norm}
With $\SLD(t)=Qt(1-Q)+(1-Q)tQ$ and $t_\delta:=\SLD^{-1}\delta$ one has, in the
eigenbasis of $Q$, $\widetilde{(t_\delta)}_{ik}=w_{ik}\widetilde\delta_{ik}$, hence
$\Tr(t_\delta\delta)=\sum_{ik}w_{ik}|\widetilde\delta_{ik}|^2=4g_Q(\delta)$ and
$V(t_\delta)=\sum_{ik}|\widetilde{(t_\delta)}_{ik}|^2q_k(1-q_i)
=\tfrac12\sum_{ik}w_{ik}^{-1}|\widetilde{(t_\delta)}_{ik}|^2=2g_Q(\delta)$.
\end{lemma}
The symmetrisation in the last step uses that $|\widetilde t_{ik}|^2$ is symmetric under
$i\leftrightarrow k$ for $t=t^*$.  Numerically ($n=9$, random $Q$ with $0<Q<1$, random
$\delta=\delta^*$): $\Tr(t\delta)/g_Q=4.000000$, $V(t)/g_Q=2.000000$,
$(\Tr t\delta)^2/(8V)/g_Q=1.000000$, and a random search over $4000$ test operators never
exceeded the Legendre value.  \textbf{Lemma~1.1 stands}, including
$b(\delta,\eta)=\frac14\Tr(t_\delta\eta)$ and
$\frac{\dd}{\dd\sigma}g_Q(\delta+\sigma\eta)|_0=\frac12\Tr(t_\delta\eta)$.

\paragraph{Adjudication of the factor $2$ (Remark~1.2).}
S10~(C4) defines $t_\delta$ by \emph{the same} equation $\SLD(t)=\delta$ (S10 eq.~(5)) and
then asserts $\Tr(t_\delta\delta)=2g_Q(\delta)$, $V(t_\delta)=\frac12g_Q(\delta)$.  By
Lemma~\ref{lem:norm} both are wrong under S10's own eq.~(5): the correct values are
$4g_Q$ and $2g_Q$.  The two errors are \emph{not} independent: they are exactly what one
gets by replacing $t$ by $t/2$ (linear in $t$, quadratic in $t$).  So S10's displayed
relations correspond to the normalisation $\SLD(t)=\frac12\delta$, i.e.\ $t_{\rm S10}
=\frac12\SLD^{-1}\delta$.  \emph{Verdict:} F1's diagnosis (``S10 has a factor $2$'') is
\textbf{correct}, and F1's own normalisation (F1 eq.~(5)) is the consistent one; but
F1's parenthetical in Remark~1.2 inverts the factor --- it reads ``$\SLD(t)=2\delta$
(i.e.\ $t\mapsto t/2$)'', and $\SLD(t)=2\delta$ means $t\mapsto2t$.  It must read
$\SLD(t)=\frac12\delta$.  \emph{Consequences:} none downstream.  The Legendre ratio
$(\Tr t\delta)^2/(8V(t))$ is homogeneous of degree $0$ in $t$, so S10 Cor.~3.3's
$\mathcal D_z\ge\Delta_z(t_c)^2/(8V(t_c))=g_Q(\delta_c)$ is correct in either
normalisation, as is the weak-duality bound of F1 Remark~3.18; and every sign
condition of \S4 is homogeneous of degree $1$ in $t_*$.  The one place where it matters is
F2's reporting convention: $\lambda_{\min}(M_*)$ ``in units of $g_Q(\dot\delta)$'' is
normalisation dependent by a factor $2$, and F2 must state which $t_*$ it uses.
\begin{verdictok}[\;\S1]
Lemma~1.1 verified analytically and numerically.  Remark~1.2: right conclusion, inverted
parenthetical (repair R1).  S10~(C4) carries a genuine internal inconsistency (repair R2,
for the S10 errata log, not for this document).
\end{verdictok}

\section{\S2: the exact chart of $\cF$ (Lemmas 2.1--2.4, Prop.~2.5)}\label{sec:chart}

\paragraph{Lemma~2.1 (the $(N,Y)$ chart), both directions.}  Recomputed: with
$N=1-XX^*$, $Y=Z-XQ_{BB}X^*$ and $R_++R_-=1_{\hh_B}$,
\[ Z-XR_+X^*=Y,\qquad 1-XR_-X^*-Z=1-XR_-X^*-XR_+X^*-Y=1-XX^*-Y=N-Y, \]
so $(X,Z)\in\cF\iff Y\ge0$ and $Y\le N$; and $Y\ge0$, $Y\le N$ force $N\ge0$.  Both
implications are equivalences of the \emph{same} two operator inequalities, no extra
hypothesis enters (in particular no boundedness of $X$ beyond $X$ bounded, and no
invertibility).  The map $(X,Z)\mapsto(X,Y)$ is affine with affine inverse
$Z=XQ_{BB}X^*+Y$.  \textbf{Stands} (exactly as written).

\paragraph{Lemma~2.2 and Remark~2.3 (the co-isometry).}
(1)--(2) are the standard polar decomposition: $X=U|X|$, $|X^*|U=U|X|U^*U=U|X|=X$, so
$X=(1-N)^{1/2}V$ with $V:=U$; $VV^*$ is the projection onto
$\overline{\ran X}=\overline{\ran XX^*}=\overline{\ran(1-N)}=\ker(1-N)^\perp
=(1-\mathbf1_{\{1\}}(N))\hh_D$.  \textbf{Stands.}

Remark~2.3's sharpness claim was verified independently.  For $X=S$ the forward shift on
$\ell^2(\mathbb N)$: $XX^*=SS^*=1-|e_0\rangle\langle e_0|$, so $N=|e_0\rangle\langle e_0|$
and $(1-N)^{1/2}=P_{\ge1}$.  Any factorisation $X=(1-N)^{1/2}V'$ forces $P_{\ge1}V'=S$,
i.e.\ $V'=S+e_0\varphi^*$ with $\varphi:=V'^*e_0$ free.  Then
$V'V'^*=SS^*+(S\varphi)e_0^*+e_0(S\varphi)^*+\|\varphi\|^2|e_0\rangle\langle e_0|$; the
$(e_j,e_0)$ entries with $j\ge1$ give $S\varphi=0$, hence $\varphi=0$ ($S$ injective),
while the $(e_0,e_0)$ entry gives $\|\varphi\|^2=1$.  Contradiction, so \emph{no}
co-isometric factor exists.  \textbf{The stated counterexample is correct.}

\begin{verdictgap}[\;Lemma~2.2(1) is sharp for the polar $V$, not for the existence of
\emph{some} co-isometric factor]
The abstract in F1 and item~1 of F1 \S7 read ``$X=(1-N)^{1/2}V$ with $V$ a co-isometry
\emph{iff} $1\notin\sigma_p(N)$''.  That is an overstatement.  Repeating the computation
above in general: $X=(1-N)^{1/2}V'$ forces $\ran(V'-V)\subseteq K:=\ker(1-N)$, and
$V'V'^*=1$ holds iff $J:=V'-V$ satisfies $JJ^*=P_K$ and $VJ^*=0$, i.e.\ iff $J^*$ is an
isometry $K\to\ker X$.  Hence
\[ \exists\ \text{co-isometric factor}\iff \dim\ker(1-N)\le\dim\ker X , \]
which for the shift ($\ker S=0$) is the document's counterexample, but which is satisfiable
with $1\in\sigma_p(N)$: take $\hh_B=\ell^2\oplus\ell^2$, $X(a,b)=Sa$, and
$V'(a,b)=Sa+e_0\langle f,b\rangle$ with $\|f\|=1$; then $N=|e_0\rangle\langle e_0|$,
$1\in\sigma_p(N)$, and $V'V'^*=1$.  This does not affect anything downstream: \S3 works in
$\|N\|<1$ where Lemma~2.2(3) applies.  \textbf{Verdict: stands with repair} (R3: state
(1) for the polar $V$, or state the dimension criterion).
\end{verdictgap}

\paragraph{Lemma~2.4 and Prop.~2.5.}  $U:=W^*V^*$ is an isometry since
$U^*U=VWW^*V^*=VV^*=1$, and $V^*=WW^*V^*=WU$; the converse puts $V:=U^*W^*$.  \textbf{Stands.}
Prop.~2.5: $\delta_{AD}=Q_{AD}-Q_{AB}X^*$ with $X^*=V^*(1-N)^{1/2}=WU(1-N)^{1/2}$ gives
the two displayed forms with $\Lambda=Q_{AB}W$; $\delta_{DD}=Q_{DD}-XQ_{BB}X^*-Y$ with
$XQ_{BB}X^*=(1-N)^{1/2}U^*(W^*Q_{BB}W)U(1-N)^{1/2}$.  Recomputed, correct, exact at finite
$s$; $\delta_{AA}=0$ because $Q_{\rm rec}$ has $Q_{AA}$ in its $AA$ slot by construction.
Remark~2.6's warning that $Q_{AD}\ne Q_{AB}$ is well taken and is the only place the
geometry enters.  \textbf{Stands.}

\section{\S3.1: the first-order defect (Lemmas 3.2--3.3, Prop.~3.4)}\label{sec:blocks}

\paragraph{Lemma~3.2 (exact feasibility at finite $s$) --- mandated check (iii).}
Recomputed from the definitions, with no expansion.  $V(s)=(V(s)^*)^*=(W_sU_s)^*=U_s^*W_s^*$,
hence
\[ V(s)V(s)^*=U_s^*W_s^*W_sU_s=U_s^*U_s=1_{\hh_D}, \qquad
   X(s)X(s)^*=(1-sN_1)^{1/2}\,1\,(1-sN_1)^{1/2}=1-sN_1 , \]
using only (i) $W_s$ \emph{restricted} as in (C3) is a unitary $\hh_D\to\hh_B$ (S10 (C5),
where $V_c=E_BW_\lambda E_D$ is unitary --- a genuine continuum feature) and (ii) $U_s$ is
an isometry of $\hh_D$ \emph{at each finite $s$}, which is Def.~3.1(1).  Then
$N=sN_1\ge0$ needs $s\|N_1\|<1$ (as stated, and it needs $N_1$ to be a bounded
\emph{operator}, see \S\ref{sec:bdry}), and $Y=Z(s)-X(s)Q_{BB}X(s)^*=sY_1$ with
$0\le sY_1\le sN_1$.  Lemma~2.1 then gives $(X(s),Z(s))\in\cF$.  The two inequalities
$0\le sY_1\le sN_1\le1$ and $X(s)X(s)^*=1-sN_1\le1$ hold \emph{exactly}, not to leading
order.  \textbf{Stands} --- this is the strongest part of the document: the upper bound is
constructive and every competitor is a genuine channel.

\paragraph{Lemma~3.3.}  $\chi=-x^2\mathbf1_{x>0}$ has $\chi(0)=\chi'(0)=0$ and $D_\chi$ is
local, so $D_\chi E_A=0=E_AD_\chi$ and $E_A\dot\delta E_A=0$.  $E_D\dot\delta E_D=0$ is
quoted from S10 Lemma~11.1 (M\"obius rigidity: on $D$, $\chi$ \emph{is} the global M\"obius
field $-x^2$, whose generator commutes with $P$); I checked S10 Lemma~11.1 at the cited
place and it says this.  \textbf{Stands} (imported, correctly cited).

\paragraph{Proposition~3.4 --- mandated check (ii): every block recomputed from scratch.}
From $\delta=Q-Q_{\rm rec}$ with $X(s)=(1-sN_1)^{1/2}V(s)$, $V(s)^*=W_sU_s$,
$U_s=1+s\zeta+o(s)$, $(1-sN_1)^{1/2}=1-\frac s2N_1+O(s^2)$, $Z=XQ_{BB}X^*+sY_1$:
\begin{align*}
 \delta_{AD}(s)&=Q_{AD}-Q_{AB}X(s)^*=Q_{AD}-\Lambda_sU_s(1-sN_1)^{1/2},\qquad \Lambda_s=Q_{AB}W_s,\\
 &=(Q_{AD}-\Lambda_s)-\Lambda_s\bigl[s\zeta-\tfrac s2N_1\bigr]+o(s)
  =s\Bigl[\dot\delta_{AD}-Q_{AD}\zeta+\tfrac12Q_{AD}N_1\Bigr]+o(s),
\end{align*}
where $Q_{AD}-\Lambda_s=(\delta_c)_{AD}=s\dot\delta_{AD}+O(s^2)$ and $\Lambda_s=Q_{AD}+O(s)$
(the $O(s)$ error multiplies an $O(s)$ factor).  This reproduces
$\Aop_{AD}=-Q_{AD}\zeta+\frac12Q_{AD}N_1$: the $\zeta$-sign is \emph{minus} because
$X^*$, not $X$, carries $U_s$, and the $N_1$-sign is \emph{plus} with the factor $\frac12$
from the square root.  For the diagonal block, with $Z_c=W_s^*Q_{BB}W_s$,
\[ X Q_{BB}X^*=\bigl(1-\tfrac s2N_1\bigr)(1+s\zeta^*)Z_c(1+s\zeta)\bigl(1-\tfrac s2N_1\bigr)+o(s)
 =Z_c+s(\zeta^*Z_c+Z_c\zeta)-\tfrac s2\{N_1,Z_c\}+o(s), \]
so $\delta_{DD}=(Q_{DD}-Z_c)-s(\zeta^*Z_c+Z_c\zeta)+\frac s2\{N_1,Z_c\}-sY_1+o(s)$; since
$Q_{DD}-Z_c=(\delta_c)_{DD}=sE_D\dot\delta E_D+O(s^2)=O(s^2)$ by Lemma~3.3 and
$Z_c=Q_{DD}+o(1)$, the first-order part is
$\Aop_{DD}=-(\zeta^*Q_{DD}+Q_{DD}\zeta)+\frac12\{N_1,Q_{DD}\}-Y_1$.  Hermiticity checks:
$\Aop_{DD}^*=\Aop_{DD}$ (the $\zeta$-term is self-adjoint as written, \emph{not} a
commutator unless $\zeta^*=-\zeta$), and
$\Aop_{DA}=\Aop_{AD}^*=-\zeta^*Q_{DA}+\frac12N_1Q_{DA}$.  Consistency with the isometric
orbit: $\zeta=-\mathcal G$, $\mathcal G^*=-\mathcal G$ gives
$\Aop_{AD}=Q_{AD}\mathcal G$ and $\Aop_{DD}=[Q_{DD},\mathcal G]$, which are the blocks of
$E_I[P,\mathcal G]E_I$ --- I verified this independently using $E_A\mathcal G=\mathcal GE_A=0$.
\textbf{All blocks and all signs confirmed.}
\begin{verdictgap}[\;Prop.~3.4: the expansion is a \emph{form} statement on $C_c^\infty$]
The proposition is stated ``as sesquilinear forms on $C_c^\infty(A)\oplus C_c^\infty(D)$'',
and in that sense it is correct.  It is however weaker than it looks: for unbounded $\zeta$
the remainder is $o(s)$ only after testing against fixed smooth compactly supported
vectors, and a boundary layer of width $O(s)$ at $\partial D$ is invisible to it.  The
entire weight of the passage from ``forms'' to ``$g_Q$'' is carried by H1, and
\S\ref{sec:bdry} shows that this is not a technicality: there are directions admitted by
Def.~3.1(1)--(2) (as forms) for which $g_Q(\delta(s))\sim s$, not $s^2$.  \textbf{Verdict:
stands as a form identity} (repair R4: say so in the statement, not only in the proof).
\end{verdictgap}

\section{\S3.2: the tangent programme and the upper bound}\label{sec:tangent}

\paragraph{Lemma~3.6 (convexity).}  The cone property was rechecked: if $\zeta_1,\zeta_2$
generate isometry families $U^1,U^2$ and $\alpha,\beta\ge0$, the product family
$s\mapsto U^1_{\alpha s}U^2_{\beta s}$ consists of isometries and has derivative
$\alpha\zeta_1+\beta\zeta_2$, so the admissible $\zeta$'s form a convex cone --- correct as
stated.  $\{N_1\ge0,0\le Y_1\le N_1\}$ is a convex cone, $\Aop$ is real-linear, so
$\mathcal C$ is a convex cone.  $g_Q\ge\frac14\|\cdot\|_2^2$ (Prop.~3.14) makes $g_Q$
strictly positive definite on $\mathfrak F$, hence $\eta\mapsto g_Q(\dot\delta+\eta)$ is
strictly convex and its kernel on $\mathcal C$ is $\{0\}$.  \textbf{Stands.}
\emph{Caveat (used in \S\ref{sec:kkt}):} a convex cone is \emph{not} a linear space.  The
admissible $\zeta$-cone is genuinely one-sided --- $\zeta=-\lambda D_\chi|_D$ ($\lambda>0$)
is admissible (inward flow), $\zeta=+\lambda D_\chi|_D$ is not (outward flow leaves $D$;
its adjoint fails dissipativity) --- and $\mathcal C$ is not closed, so the minimiser need
not be attained.  Neither point is flagged in the document.

\paragraph{H1 (Def.~3.7) and Theorem~3.9.}  The proof is the triangle inequality for the
Hilbert seminorm $g_Q^{1/2}$ plus H1, applied to a family that Lemma~3.2 has already shown
to be \emph{exactly} feasible; squaring and dividing by $s^2$ gives \eqref{eq:upperref}:
\begin{equation}\label{eq:upperref}
 \limsup_{s\downarrow0}s^{-2}\min_\cF g_Q\ \le\ g_Q\bigl(\dot\delta+\Aop(\zeta,N_1,Y_1)\bigr).
\end{equation}
\textbf{Stands, given H1 for the datum used.}  Two honest points which the document makes
and which I confirm: H1 is \emph{not} implied by Prop.~3.4 (the $w_{ik}$ grow like
$e^{\min(|\kappa_i|,|\kappa_k|)}$, so $g_Q$ is strictly stronger than $\|\cdot\|_2$), and
H1 is known on the isometric orbit (S10 Lemma~6.1$\langle1\rangle3$--$\langle1\rangle4$
carries out exactly that estimate --- checked at the cited place).  The statement
``$\limsup\le T$ if H1 holds for a minimising sequence'' is correct but should be read
with \S\ref{sec:bdry}: which data are admitted in the infimum $T$ is exactly what
Thm.~3.11 gets wrong, and for some of them H1 \emph{fails}.

\paragraph{Corollary~3.10.}  $\Aop(-\mathcal G,0,0)=E_I[P,\mathcal G]E_I$ was verified
blockwise above; with S10 Prop.~4.3 ($g_Q(E_I[P,\mathcal Z]E_I)=\frac12\|\Pi P^\perp\mathcal ZP\|_2^2$
for anti-self-adjoint $\mathcal Z$) this gives
$g_Q(\delta^{(1)})=\frac12\|u+v_{\mathcal G}\|_2^2$ and
$\inf_{\mathcal G}=\frac12\dist(u,\mathcal Y)^2=(1-\theta)g_Q(\dot\delta)$ with $\theta$
of S10 Thm.~6.2 (sup over \emph{bounded anti-self-adjoint} $\mathcal G=E_D\mathcal GE_D$ ---
I checked the class at the cited place).  Hence $T\le(1-\theta)g_Q(\dot\delta)$.
\textbf{Stands.}  Note that the infimum is over a class whose $\{v_{\mathcal G}\}$ need not
be closed: $\delta_*$ exists as the $g_Q$-orthogonal projection residual of $-\dot\delta$
onto $\overline{\mathcal Y}$, but a minimising \emph{operator} $\mathcal G_*$ need not
exist.  \S\ref{sec:kkt} takes this up.

\section{\S3.3: boundary-moving generators (Thm.~3.11) --- the main failure}\label{sec:bdry}

This is mandated check~(iv).  I take the four items in turn.

\paragraph{(0) The definition of $\sigma$ is not well posed as written.}
$\sigma:=\frac12(\zeta+\zeta^*)$ is declared ``on $\operatorname{dom}\zeta^*$'', but $\zeta$
is \emph{not} defined there (in the document's own example
$\operatorname{dom}\zeta=H^1_0(D)\subsetneq H^1(D)=\operatorname{dom}\zeta^*$), and on the
common domain $\operatorname{dom}\zeta^{**}\cap\operatorname{dom}\zeta^*$ one has
$\sigma=0$ identically.  The object actually used in \eqref{eq:sigmarefer} is the
\emph{boundary form}
\begin{equation}\label{eq:sigmarefer}
 \langle\sigma g,g\rangle:=\Re\langle\zeta^*g,g\rangle ,\qquad g\in\operatorname{dom}\zeta^* ,
\end{equation}
which is what Remark~3.12 evaluates.  Repair R5: define $\sigma$ by \eqref{eq:sigmarefer}.
I verified Remark~3.12 independently: for $\zeta=-D_\psi$, $\zeta^*=D_\psi$ on $H^1(D)$,
$2\Re\langle D_\psi g,g\rangle=\int_D(\psi|g|^2)'=[\psi|g|^2]_0^1$, so
$-2\sigma=|\psi(1)|\delta_1\otimes\delta_1+\psi(0)\delta_0\otimes\delta_0\ge0$ for inward
$\psi$.  I also checked the domain claim implicit there: if $\psi(0)>0$ the isometry
$U_s$ leaves a hole of width $\sim s\psi(0)$ at the left endpoint, so
$\|U_sf-f\|_2^2\gtrsim s\psi(0)|f(0)|^2$ and $\operatorname{dom}\zeta=H^1_0(D)$ indeed.

\paragraph{(1) $\sigma\le0$: true, but the proof is not available in the stated generality.}
The proof says ``$U_s^*$ is a contraction semigroup with generator $\zeta^*$'' and invokes
Lumer--Phillips.  Def.~3.1(1) does \emph{not} require $(U_s)$ to be a semigroup, only a
differentiable family, and Lumer--Phillips is a statement about semigroup generators.  From
a bare family one gets $\Re\langle\zeta f,f\rangle=0$ on $\operatorname{dom}\zeta$
(differentiate $\|U_sf\|^2=\|f\|^2$) but \emph{not} $\Re\langle\zeta^*g,g\rangle\le0$ for
$g\notin\operatorname{dom}\zeta$: that needs $s\mapsto U_s^*g$ to be differentiable, which
does not follow.  Repair R6: add ``semigroup'' to Def.~3.1(1) (every example in the
document, $e^{s\mathcal G}$ and $W_{\varphi_s}$ for a flow $\varphi$, is one), or assume
weak differentiability of the adjoint family.  Under either, $\sigma\le0$ is correct, and
it is \emph{essential}: it is the only thing that excludes the outward field of~(3).
\textbf{Stands with repair.}

\paragraph{(2) The absorption identity: algebra correct, content empty.}
Substituting $\zeta=\mathcal G+\sigma$, $\zeta^*=-\mathcal G+\sigma$ into the blocks gives,
as I recomputed,
$\Aop_{DD}=[\mathcal G,Q_{DD}]+\frac12\{N_1-2\sigma,Q_{DD}\}-Y_1$ and
$\Aop_{AD}=-Q_{AD}\mathcal G+\frac12Q_{AD}(N_1-2\sigma)$, i.e.\
$\Aop(\zeta,N_1,Y_1)=\Aop(\mathcal G,N_1-2\sigma,Y_1)$.  But Prop.~3.4 lives on
$C_c^\infty(A)\oplus C_c^\infty(D)$, where $\sigma$ \emph{vanishes identically}
($g(0)=g(1)=0$); there the identity reads $\Aop(\zeta,N_1,Y_1)=\Aop(\zeta,N_1,Y_1)$ and
$\mathcal G$ is the same differential expression as $\zeta$.  On any larger domain the
left-hand side is undefined.  So the identity is a tautology and does \emph{not} show that
a boundary-moving direction ``is'' an $N_1$-direction.  (It \emph{is} the right statement
for \emph{contraction} families: if $U_s$ is only a contraction, $1-U_s^*U_s=-2s\sigma+o(s)
\ge0$ is an honest loss of isometry and the absorption is exact.  That class is not
Def.~3.1(1).)  \textbf{Stands as algebra; the interpretation fails.}

\begin{verdictbreak}[\;Theorem~3.11(3) is false as stated, and fatally so]
Step $\langle1\rangle4$ asserts: ``conversely $\mathcal G$ anti-symmetric with
$\mathcal G=E_D\mathcal GE_D$ generates the isometry (indeed unitary) family
$e^{s\mathcal G}$, so the two cones coincide''.  Counterexample:
$\mathcal G:=+\lambda D_\chi|_D$, $\lambda>0$, is anti-symmetric on $C_c^\infty(D)$
($D_\chi^*=-D_\chi$) and supported in $D$, but $e^{s\mathcal G}$ is the half-density
push-forward along the flow of $+\lambda\chi$, i.e.\ $x\mapsto x/(1-s\lambda x)$, which
maps $D=(0,1)$ \emph{onto} $(0,\frac1{1-s\lambda})\supsetneq D$: it is not an operator on
$\hh_D$, and its boundary form is $\Re\langle-\lambda D_\chi g,g\rangle
=+\frac\lambda2|g(1)|^2>0$, violating item~(1).  It is therefore \emph{not} admissible ---
which the document itself says in Remark~3.13 (``the opposite sign $\lambda<0$
(decompression) is inadmissible: it points out of $D$'').  The two statements contradict
each other, and the wrong one is load-bearing: with $\zeta=+\lambda D_\chi$ one gets
$\Aop_{AD}=-\lambda Q_{AD}D_\chi=-\lambda\dot\delta_{AD}$, $\Aop_{DD}=\lambda[Q_{DD},D_\chi]=0$,
hence $\delta^{(1)}=(1-\lambda)\dot\delta$ and $g_Q(\delta^{(1)})=(1-\lambda)^2g_Q(\dot\delta)
\to0$ as $\lambda\to1$: the enlarged cone has tangent value $T=0$, i.e.\
$\kappa_{\rm opt}=0$.  (No contradiction with reality: for that family $U_s$ is a
contraction, not an isometry, the true channel is $X^*=$ the restriction $\hh_D\to\hh_B$,
its defect block is $Q_{AC}$ with $\|Q_{AC}\|_2^2=O(s)$, and $g_Q(\delta(s))=O(s)\ne O(s^2)$:
H1 fails for it, catastrophically.)
\emph{What is true:} $\mathcal C\subseteq\{\Aop(\mathcal G,N_1',Y_1)\}$ with $\mathcal G$
anti-symmetric \emph{and an isometry generator} and $N_1'\ge0$ a form.  That class is a
one-sided cone, not a linear space.  Repair R7.
\end{verdictbreak}

\paragraph{(3$'$) The ``closure'' claim of F1 \S7 item 4 is also wrong in the other
direction.}  The boundary elements are not in the $g_Q$-closure of the operator directions:
for $N_1^{(n)}:=\lambda n\mathbf1_{(1-1/n,1)}\rightharpoonup\lambda\delta_1\otimes\delta_1$
one has $\|Q_{AD}N_1^{(n)}\|_2^2\approx\lambda^2n/(4\pi^2)\to\infty$, so
$\Aop(0,N_1^{(n)},0)$ diverges in $g_Q$.  The singular directions sit at infinity, not in
the closure.

\paragraph{(4) Bounded $\zeta\Rightarrow\sigma=0$.}  Correct and independent of
Lumer--Phillips: $\operatorname{dom}\zeta=\hh_D$, $\Re\langle\zeta f,f\rangle=0$ for all
$f$, and a bounded symmetric operator is self-adjoint, so $\zeta^*=-\zeta$.
\textbf{Stands.}  This is the one part of \S3.3 that survives intact, and it is worth
stating as the headline of the section: \emph{bounded} isometry families give nothing
beyond S10's class; \emph{unbounded} ones are an extra one-sided cone that must be tested
separately (\S\ref{sec:kkt}).  Remark~3.13 (over-compression) was verified in full:
$\delta^{(1)}=(1+\lambda)\dot\delta$, one-sided derivative $2\lambda g_Q(\dot\delta)>0$.

\section{\S3.4: the lower bound, H2, and the honesty of Remark 3.15}\label{sec:lower}

\paragraph{Prop.~3.14.}  $a(1-b)+b(1-a)=a+b-2ab$ is affine in each variable, so on
$[0,1]^2$ it is maximised at a vertex, with values $0,1,1,0$; hence $w_{ik}\ge1$ and
$g_Q\ge\frac14\|\cdot\|_2^2$.  (In fact $w_{ik}=1+\cosh\frac{\kappa_i+\kappa_k}2/
\cosh\frac{\kappa_i-\kappa_k}2>1$ strictly --- verified numerically.)  The block bounds
\eqref{eq:HSref} follow since each block norm is $\le\|\delta\|_2$.  \textbf{Stands.}
\begin{equation}\label{eq:HSref}
 \|Q_{AD}-Q_{AB}X^*\|_2\le2\sqrt Cs,\qquad \|Q_{DD}-XQ_{BB}X^*-Y\|_2\le2\sqrt Cs .
\end{equation}

\paragraph{Remark~3.15 --- mandated check (viii), the essential-spectrum obstruction.}
The remark is \textbf{honest and correct}, and is the best part of the lower-bound
discussion.  Three points, all of which I confirm: (i) to turn \eqref{eq:HSref} into
$\|X-X_c\|=O(s)$ one must invert $Q_{AB}$ from the left; (ii) $Q_{AB}=E_APE_B$ is the
off-diagonal block of $Q=E_IPE_I$ for two intervals sharing the endpoint $0$, whose
modular spectrum fills $(0,1)$, so the singular values of $Q_{AB}$ accumulate at $0$ and no
bounded left inverse exists --- the invisible directions of Phase~4 are exactly the
unresolved ones; (iii) the $DD$ estimate controls only the combination $XQ_{BB}X^*+Y$, so
$N$ and $Y$ are not separately localised.  Nothing is claimed that is not proved.

\paragraph{H2 (Def.~3.16) and Theorem~3.17.}  The $\liminf$ step is the triangle
inequality again and is correct.  But H2 as formulated assumes (a) that a minimiser
\emph{exists} at each small $s$, (b) that it is \emph{of the form} \eqref{eq:scaledref}
with a first-order datum, (c) uniformly bounded $\|N_1^{(s)}\|+\|Y_1^{(s)}\|$ and
(d) a uniform $o(s^2)$ remainder.  That is very close to assuming the conclusion: (b)+(d)
already say that the minimiser is a differentiable curve through the compression with the
tangent structure of \S3.1.  The document labels it a hypothesis and repeats the label in
the abstract, in the verdict box after Thm.~4.4 and in \S7, so there is no dishonesty ---
but it should be said explicitly that H2 is not a regularity assumption, it is the lower
bound in disguise.  Repair R8.
\begin{equation}\label{eq:scaledref}
 X(s)=(1-sN_1)^{1/2}V(s),\quad V(s)^*=W_sU_s,\quad Z(s)=X(s)Q_{BB}X(s)^*+sY_1 .
\end{equation}

\paragraph{Remark~3.18 (the unconditional route).}  The displayed formula for $\sup_\cF L$
was rederived: with S10's $V=X^*$, $L=2\Re\Tr(t_{DA}Q_{AB}V)+\Tr(t_{DD}V^*R_+V)+\Tr(t_{DD}Y)$
and $\sup_{0\le Y\le1-V^*V}\Tr(t_{DD}Y)=\Tr[((1-V^*V)^{1/2}t_{DD}(1-V^*V)^{1/2})_+]$ by the
Douglas argument of Thm.~4.4$\langle1\rangle1$.  Correct.  One caveat to add: the weak
duality bound $\min_\cF g_Q\ge\Delta_z(t)^2/(8V(t))$ requires $\Delta_z(t)\ge0$ (otherwise
$\inf_\cF(\Tr t\delta)^2$ is not bounded below by $\Delta_z(t)^2$); for $t=t_c$ this is
S10 Cor.~3.3, but for a general $t$ it must be checked.  Repair R9.

\section{\S4: the KKT criterion (Lemma 4.2, Prop.~4.3, Thm.~4.4, Cor.~4.5, Prop.~4.6)}\label{sec:kkt}

\paragraph{Lemma~4.2, derivative formula --- mandated check (v), first part.}
Recomputed from scratch.  $\frac{\dd}{\dd\sigma}g_Q(\delta_*+\sigma\eta)|_0=\frac12\Tr(t_*\eta)$
(Lemma~\ref{lem:norm}), and with $\Aop_{AA}=0$,
\begin{align*}
 \Tr(t_*\Aop)&=\Tr\bigl(t_{*,AD}\tfrac12N_1Q_{DA}\bigr)+\Tr\bigl(t_{*,DA}\tfrac12Q_{AD}N_1\bigr)
  +\Tr\bigl(t_{*,DD}(\tfrac12\{N_1,Q_{DD}\}-Y_1)\bigr)\\
 &=\Tr\bigl(N_1\Herm(t_{*,DA}Q_{AD})\bigr)+\Tr\bigl(N_1\Herm(t_{*,DD}Q_{DD})\bigr)-\Tr(Y_1\tau_*)
  =\Tr(N_1M_*)-\Tr(Y_1\tau_*),
\end{align*}
using $Q_{DA}t_{*,AD}=(t_{*,DA}Q_{AD})^*$ and $(t_*Q)_{DD}=t_{*,DA}Q_{AD}+t_{*,DD}Q_{DD}$
(valid because $E_A+E_D=1_{\hh_I}$).  So the derivative is
$\frac12[\Tr(N_1M_*)-\Tr(Y_1\tau_*)]$ with $M_*=E_D\Herm(t_*Q)E_D$, $\tau_*=E_Dt_*E_D$,
exactly as claimed; the \emph{normalisation} is F1's (Lemma~\ref{lem:norm}), i.e.\ a factor
$2$ larger than it would be in S10's convention --- see \S\ref{sec:norm}.  Numerically
(random $Q$, $\delta_*$, $N_1\ge0$, $Y_1=0.3N_1$, $n=9$) the symmetric difference quotient
of $g_Q$ agrees with the formula to $10^{-9}$ relative, and the two definitions of $M_*$
(blockwise and $\frac12E_D(t_*Q+Qt_*)E_D$) agree to machine precision.  \textbf{Stands.}

\begin{verdictbreak}[\;The last clause of Lemma~4.2 is false]
``Along the isometric directions $\Aop(\mathcal G,0,0)$ ($\mathcal G$ anti-symmetric) the
derivative vanishes'' --- take $\mathcal G=-\lambda D_\chi|_D$, $\lambda>0$, which is
anti-symmetric, supported in $D$, and \emph{admissible} (it generates the inward
compression semigroup; this is the document's own Remark~3.13).  Then
$\Aop(\mathcal G,0,0)=\lambda\dot\delta$ and the derivative is
\[ \tfrac12\Tr(t_*\lambda\dot\delta)=2\lambda\,b(\delta_*,\dot\delta)
 =2\lambda\bigl[g_Q(\delta_*)-b(\delta_*,[Q,\mathcal G_*])\bigr]=2\lambda(1-\theta)g_Q(\dot\delta)>0 , \]
using the normal equation $b(\delta_*,[Q,\mathcal G])=0$ for $\mathcal G$ in the class over
which $\delta_*$ is optimal.  The justification given ($\mathcal G\mapsto\Aop(\mathcal G,0,0)$
takes values in a \emph{linear space} $L$ on which $\delta_*$ minimises) fails precisely
because the admissible anti-symmetric generators form a one-sided cone, not a linear space
(\S\ref{sec:bdry}).  Repair R10: state the clause for $\mathcal G$ in the linear class of
S10 (bounded anti-self-adjoint, or the closed subspace $\overline{\mathcal Y}$), and treat
the unbounded ones as one-sided directions with their own sign condition
$-b(\delta_*,[Q,\zeta])\ge0$.  Note that this extra condition \emph{holds} for the
over-compression direction (the display above is $>0$), so nothing is broken numerically;
what is broken is the claim that Thm.~4.4(a) is \emph{complete}.
\end{verdictbreak}

\paragraph{Prop.~4.3.}  $\Tr(t_*[\mathcal G,Q])=\Tr(E_D[Q,t_*]E_D\,\mathcal G)$;
$C:=E_D[Q,t_*]E_D$ is anti-Hermitian and $\Tr(C\mathcal G)=0$ for all anti-Hermitian
$\mathcal G$ forces $C=0$ (take $\mathcal G=C$: $\Tr(C^2)=-\|C\|_2^2$), whence
$(t_*Q)_{DD}=(t_*Q)_{DD}^*=M_*$.  Correct, with two caveats: the test direction
$\mathcal G=C$ must itself be admissible (bounded, supported in $D$) --- true if $C$ is
bounded, otherwise a density argument is needed --- and ``stationary along the isometric
directions'' must again mean the linear class.  \textbf{Stands with repair (R10).}

\paragraph{Theorem~4.4, (a)$\Leftrightarrow$(c) --- mandated check (v), second part.}
Both directions were rederived and are airtight.  (a)$\Rightarrow$(c): for $0\le Y_1\le N_1$,
$\Tr(Y_1\tau_*)\le\Tr(Y_1M_*)\le\Tr(N_1M_*)$ (traces of products of positive operators are
$\ge0$), and the middle term is maximised by Douglas: $Y_1\le N_1\Rightarrow
Y_1^{1/2}=N_1^{1/2}K$, $\|K\|\le1$, so $Y_1=N_1^{1/2}CN_1^{1/2}$ with $C=KK^*\in[0,1]$ and
$\max_C\Tr(C\,N_1^{1/2}\tau_*N_1^{1/2})=\Tr[(N_1^{1/2}\tau_*N_1^{1/2})_+]$ at
$C=\mathbf1_{(0,\infty)}$.  (c)$\Rightarrow$(a): rank-one $N_1=|\psi\rangle\langle\psi|$ is
a projection, $N_1^{1/2}\tau_*N_1^{1/2}=\langle\psi,\tau_*\psi\rangle N_1$, and (c) reads
$\langle\psi,M_*\psi\rangle\ge\max\{0,\langle\psi,\tau_*\psi\rangle\}$.  (b)$\Rightarrow$(a)
is immediate from $(\tau_*)_+\ge\max\{0,\tau_*\}$, and is strictly stronger because the
operator order is not a lattice order.
\textbf{Counterexample search (mandated):} see \S\ref{sec:numcheck}; none exists, as the
proof predicts.


\subsection{Numerical counterexample search for (a)$\Leftrightarrow$(c)}\label{sec:numcheck}
$40\,000$ random pairs $(M_*,\tau_*)$ of Hermitian $2\times2$ and $3\times3$ matrices with
random $N_1\ge0$ (ranks $1$ and $m$), all generically non-commuting (typical
$\|[M_*,\tau_*]\|\sim6$, $\|[M_*,N_1]\|\sim2$): \textbf{zero} cases where (a) holds and (c)
fails, and in all $39\,278$ cases where (a) failed the rank-one witness of
Thm.~4.4$\langle1\rangle3$ (the bottom eigenvector of $M_*$ or of $M_*-\tau_*$) detected it.
A second, independent run with a different seed and sample size reproduced this
($0/20\,000$, $2938/2938$).  The Douglas step was checked separately: over $60$ random
$(N_1,\tau_*)$ with $800$ random admissible $Y_1$ each, no $Y_1$ ever beat
$\Tr[(N_1^{1/2}\tau_*N_1^{1/2})_+]$ (worst excess $0.000\rm e{+}00$).
\textbf{No counterexample exists, in agreement with the proof.}

\subsection{Robustness: is $M_*\ge0$ meaningful in the continuum?}\label{sec:robust}
This is where I part company with the document's presentation, and it bears directly on
F2's negative result.  $t_*=\SLD^{-1}\delta_*$ has
$\|t_*\|_2^2=\sum_{ik}w_{ik}^2|\widetilde{(\delta_*)}_{ik}|^2$ while
$g_Q(\delta_*)=\frac14\sum_{ik}w_{ik}|\widetilde{(\delta_*)}_{ik}|^2<\infty$; since
$w_{ik}=1+\cosh\frac{\kappa_i+\kappa_k}2/\cosh\frac{\kappa_i-\kappa_k}2$ is unbounded,
$\delta_*\in\mathfrak F$ does \emph{not} imply $t_*\in\Stwo$, let alone trace class.  So
$M_*$ and $\tau_*$ are in general \emph{quadratic forms}, not operators, $\Tr(N_1M_*)$ is
not a trace, and ``$M_*\ge0$'' has no spectral meaning without a domain.  The repair is
available inside the document's own framework and costs nothing: define
\begin{equation}\label{eq:formdef}
 \langle\psi,M_*\psi\rangle:=4\,b\bigl(\delta_*,\Aop(0,|\psi\rangle\langle\psi|,0)\bigr),\qquad
 \langle\psi,\tau_*\psi\rangle:=-4\,b\bigl(\delta_*,\Aop(0,0,|\psi\rangle\langle\psi|)\bigr),
\end{equation}
on $\mathfrak D:=\{\psi:\Aop(0,|\psi\rangle\langle\psi|,0)\in\mathfrak F\}$.  These are
finite for every $\psi\in\mathfrak D$ by Cauchy--Schwarz for $b$, they agree with
Def.~4.1 whenever $t_*$ is trace class, and (a) becomes the pair of \emph{form}
inequalities on $\mathfrak D$.  With that reading Thm.~4.4 is a correct statement;
$\lambda_{\min}(M_*)$ of \S5 item~4 is then $\inf\{\langle\psi,M_*\psi\rangle:\psi\in\mathfrak D,\|\psi\|=1\}$
and may be $-\infty$.  Repair R11.

\begin{verdictgap}[\;Corollary 4.5: the gain can vanish, so strictness needs care]
The algebra is right: $g_Q(\delta_*+\sigma d)=g_Q(\delta_*)+2\sigma\Delta+\sigma^2g_Q(d)$ is
minimised at $\sigma_{\rm opt}=-\Delta/g_Q(d)>0$ with value $g_Q(\delta_*)-\Delta^2/g_Q(d)$,
and Lemma~3.2 makes the improved datum a genuine channel.  But the conclusion
``$\kappa_{\rm opt}<1-\theta$ strictly'' needs \emph{one} admissible $\psi$ with
$\Delta<0$ \emph{and} $g_Q(d)<\infty$.  If the form \eqref{eq:formdef} is unbounded below
only along a sequence $\psi_n$ with $\Delta_n^2/g_Q(d_n)\to0$, then
$T=(1-\theta)g_Q(\dot\delta)$ still, with the infimum not attained, and
$\kappa_{\rm opt}=1-\theta$.  This is exactly the configuration F2 reports (T2:
$\lambda_{\min}(M_*)<0$ and $\lambda_{\min}(M_*-\tau_*)<0$ at every $L$, but achievable
rank-one gain $\sim10^{-10}g_Q(\dot\delta)$, and a violation that scales with the weight
cap, i.e.\ with the modular-UV cutoff that regularises $w$).  \textbf{F2's numbers do not
establish $\kappa_{\rm opt}<1-\theta$}; they are equally consistent with a form that is
unbounded below in the modular ultraviolet with vanishing achievable gain.  What would
settle it: $\sup_\psi\Delta^2/g_Q(d)$ as a function of $L$ \emph{and} of the weight cap
$\kappa_c$, with a limit bounded away from $0$.  A second caveat for F2: the circle model
violates M\"obius rigidity at $O(1/L)$, and rigidity ($E_D\dot\delta E_D=0$, Lemma~3.3) is
used in Prop.~3.4$\langle1\rangle5$; the document's own free check
$\mathrm{AntiHerm}((t_*Q)_{DD})=0$ (\S5 item~6) is the right arbiter and must be reported
alongside the two eigenvalues.
\end{verdictgap}

\paragraph{Prop.~4.6 and Remark~4.7 --- mandated check (vi).}
$\Sigma_2=t_{DA}Q_{AB}V_c+t_{DD}V_c^*R_+V_c$ and $V_c=W\to1$, $Q_{AB}\to Q_{AD}$,
$V_c^*R_+V_c=Z_c\to Q_{DD}$ strongly, so $\Sigma_2\to(tQ)_{DD}$ and
$\Sigma_1=\Sigma_2-t_{DD}\to(tQ)_{DD}-t_{DD}$: correct (I checked S10 Def.~3.1 and
Thm.~3.2 at the cited places).  The identification $\Sigma_2^{(0)}=M$ requires (S), as
stated.  Remark~4.7's ``verbatim'' claim was verified in detail: with $\gamma\ge0$ on
$\hh_B$, $N_1:=X_c\gamma X_c^*$, $\sigma_2=X_c^*\Sigma_2X_c$, $\hat t=X_c^*t_{DD}X_c$ and
$X_c$ \emph{unitary}, one has $\Tr(\sigma_2\gamma)=\Tr(\Sigma_2N_1)$ and
$X_c(\gamma^{1/2}\hat t\gamma^{1/2})X_c^*=N_1^{1/2}t_{DD}N_1^{1/2}$, so S10
Thm.~3.2$\langle1\rangle8$ is literally criterion~(c) at the compression.  \textbf{Stands},
with three caveats to record: the limits are strong (only ``$\ge0$ for all $s$'' passes to
the limit); (S) is needed for $\Sigma_2^{(0)}=M$; and S10's certificate assumes $t$ trace
class, which $t_*$ need not be (\S\ref{sec:robust}).

\section{\S6: the modular frame (Lemma 6.1, Prop.~6.2) and the superseded $\theta$-diagnosis}\label{sec:theta}

\paragraph{Lemma~6.1.}  $D_\chi$ is local and $\chi\in C^{1,1}$ vanishes to second order at
$0$, so $D_\chi f$ is supported in $\operatorname{supp}f$; hence
$(1-E_I)D_\chi E_I=0=E_ID_\chi(1-E_I)$ and, inserting $1=E_I+(1-E_I)$ between $P$ and
$D_\chi$ in both terms of $E_I[P,D_\chi]E_I$, $\dot\delta=Qd-dQ=[Q,d]$ with $d=E_ID_\chi E_I$.
Recomputed, correct.  \textbf{Stands.}

\paragraph{Proposition~6.2 --- mandated check (vii).}  All three ingredients were
recomputed from scratch.
\emph{(a) The change of variables.}  $y=-\log(1-x)$ maps $I=(-\infty,1)$ onto $\mathbb R$
with $A\mapsto\{y<0\}$, $D\mapsto\{y>0\}$; $\dd y/\dd x=1/(1-x)$.
\emph{(b) The push-forward of the field.}  A vector field $\chi\partial_x$ becomes
$w\partial_y$ with $w=\chi\,\dd y/\dd x=\chi/(1-x)$, and the half-density push-forward
intertwines $D_\chi$ with $D_w$ (both are the generators of the respective flows on
half-densities).  With the \emph{same} $\chi$ as S10~(N5) --- $\chi=-x^2\mathbf1_{x>0}$,
the generator of $x\mapsto x/(1+sx)$, so that $B=(0,\frac1{1+s})$; checked at the cited
place --- and $x=1-e^{-y}$:
\[ w=-\frac{x^2}{1-x}=-(1-e^{-y})^2e^{y}=-(e^{y/2}-e^{-y/2})^2=-4\sinh^2\tfrac y2 , \]
supported in $y>0$, $\sim-y^2$ at $0$ (so $C^{1,1}$) and $\sim-e^{y}$ at $+\infty$.
\textbf{Confirmed, including the sign.}
\emph{(c) The kernel.}  $\widetilde{[Q,X]}_{ik}=(q_i-q_k)\widetilde X_{ik}$ and, with
$a=\kappa_i$, $b=\kappa_k$, $\mathsf s=\frac{a+b}2$, $\mathsf t=\frac{a-b}2$,
$(1+e^a)(1+e^b)=2e^{\mathsf s}(\cosh\mathsf s+\cosh\mathsf t)$, $e^b-e^a=-2e^{\mathsf s}\sinh\mathsf t$,
$e^a+e^b=2e^{\mathsf s}\cosh\mathsf t$, so
$q_i-q_k=-\sinh\mathsf t/(\cosh\mathsf s+\cosh\mathsf t)$,
$w_{ik}^{-1}=\cosh\mathsf t/(\cosh\mathsf s+\cosh\mathsf t)$ and
$w_{ik}(q_i-q_k)^2=\sinh^2\mathsf t/[\cosh\mathsf t(\cosh\mathsf s+\cosh\mathsf t)]$
(verified numerically at random $(a,b)$ to $10^{-15}$).  With $\kappa=-2\pi p$ this is the
displayed $K(p,p')$; both arguments enter through even functions, so the sign convention in
\eqref{eq:Qmultref} is indeed immaterial --- I checked that $K$ is invariant under
$(p,p')\mapsto(-p,-p')$.  \textbf{Prop.~6.2 stands.}  Two things to state that are
currently implicit: the Plancherel normalisation in which
$\Tr(X^*X)=\iint|\widehat X|^2$, and that $g_Q$ is frame-independent because it is built
from the spectral data of $Q$ alone (the document says this in $\langle1\rangle1$).
\begin{equation}\label{eq:Qmultref}
 Q=(1+e^{-2\pi\mathsf p})^{-1},\qquad \mathsf p=-i\partial_y,\qquad \kappa=-2\pi p .
\end{equation}

\begin{verdictbreak}[\;Remark 6.4 and the closed-form verdict box are superseded by F3]
The document's structural diagnosis is: ``the discrepancy is a statement about how much
weight the region $y\gg1$ carries \dots\ the two discretisations are not expected to agree
before $h$ resolves $e^{-y}$-scales.  \textbf{Status: unresolved}; $\theta$ is known only
as $0.30\lesssim\theta\lesssim0.40$.''  Track F3 (\texttt{numerics/optimality\_all/f3\_t0.py},
\texttt{F3\_RESULTS.md}) has since refuted both halves: the modular Galerkin frame and a
direct discretisation of \emph{this document's own} Prop.~6.2 frame agree to $0.001$ at
matched $h$ and extrapolate to $\theta=0.384\pm0.003$ (measured order $h^{0.77}$, five mesh families; quoted as $0.385\pm0.005$ in the hand-off); the modular window $\Lambda$ is
irrelevant; the residual $h$-dependence is a \emph{bulk} effect, the endpoint piece being
saturated already at coarse $h$ --- i.e.\ the opposite of ``governed by the region
$y\gg1$''.  S10's $0.300\pm0.005$ is not a competing estimate but a \emph{taper-suppressed
lower bound}: relaxing the circle model's smoothstep ultraviolet taper moves $\theta$ at
$L=256$ from $0.2896$ to $0.36$ and beyond.  Required repairs (R12): replace the window
$0.30$--$0.40$ by $\theta=0.384\pm0.003$; replace the diagnosis of Remark~6.4; record that
the three circulated candidates are now \emph{excluded} ($3/\pi^2=0.30396$, $1/3$,
$1-2/\pi=0.36338$ lie $4\sigma$ to $26\sigma$ below the reconciled value, the closest being $1-2/\pi$);
and update the two numerical asides ``$1-\theta\ge0.66$ in both models'' $\to$
$1-\theta=0.616\pm0.003$, and $\kappa_{\rm opt}=1-\theta=0.616\pm0.003$ if the criterion
holds.  \emph{Nothing else in \S6 depends on them}: Lemma~6.1, Prop.~6.2 and Remark~6.3 are
pure algebra plus S10 Cor.~11.2, all unaffected, and the document's refusal to select a
closed form was and remains the correct stance.  The one sentence of Remark~6.4 that
survives is $K(p,p')\to1$ as $|p-p'|\to\infty$ at fixed $p+p'$ (verified); note also
$K\to0$ as $|p+p'|\to\infty$ at fixed $p-p'$, which is the direction the circle model's
taper suppresses --- consistent with F3's finding that the taper biases $\theta$ downwards.
\end{verdictbreak}

\section{Open-hypothesis bookkeeping (H1, H2)}\label{sec:hyp}

\paragraph{Is anything used that is neither proved nor declared?}  Mandated check~(viii).
Six items, in decreasing order of seriousness.
\begin{enumerate}[leftmargin=1.8em,itemsep=1pt]
\item \emph{Semigroup property of $(U_s)$} --- used through Lumer--Phillips in
 Thm.~3.11(1), not in Def.~3.1(1).  Undeclared.  (R6)
\item \emph{Attainment of the orbit optimum} --- \S4 posits ``$\zeta_*=-\mathcal G_*$ with
 $\mathcal G_*$ the minimiser of Cor.~3.10'', but Cor.~3.10 gives an infimum over a class
 whose image $\{v_{\mathcal G}\}$ is not claimed closed.  Undeclared.  The repair is free:
 define $\delta_*$ as the $g_Q$-orthogonal residual of $-\dot\delta$ on
 $\overline{\mathcal Y}$, which always exists, and $t_*:=\SLD^{-1}\delta_*$; the normal
 equation $b(\delta_*,[Q,\mathcal G])=0$ then holds by construction.  (R15)
\item \emph{$t_*$ trace class / $M_*,\tau_*$ operators} --- assumed throughout \S4--\S5
 (``both self-adjoint on $\hh_D$'', $\Tr(N_1M_*)$, $\lambda_{\min}$), but
 $\delta_*\in\mathfrak F$ does not give $t_*\in\Stwo$ (\S\ref{sec:robust}).  Undeclared.  (R11)
\item \emph{Linearity of the isometric direction space} --- used in Lemma~4.2$\langle1\rangle6$
 and Thm.~4.4$\langle1\rangle4$; false for the admissible class (\S\ref{sec:bdry}).  (R10)
\item \emph{$N_1$ a bounded operator} --- Def.~3.1(2) says ``quadratic forms'', but
 Def.~3.2 needs $(1-sN_1)^{1/2}$ and Lemma~3.2 needs $s\|N_1\|<1$.  Internal
 inconsistency, undeclared.  (R13)
\item \emph{$\Delta_z(t)\ge0$} in the weak-duality remark 3.18.  (R9)
\end{enumerate}
H1 and H2 themselves are declared, repeated in the abstract, in the verdict box after
Thm.~4.4 and in \S7, and their roles (upper resp.\ lower bound) are stated correctly; the
document never claims a result that secretly uses them.  \textbf{No circularity with S10 or
E1 was found}: the upper bound (Thm.~3.9) uses only Lemma~3.2, H1 and the triangle
inequality; Cor.~3.10 imports $\theta$ from S10 Thm.~6.2, which is proved independently of
anything here; Prop.~4.6 compares with S10 Thm.~3.2 but does not use it; E1 Cor.~7.5 is not
used in any proof.

\section{Verdict table}\label{sec:table}
\newcommand{\vS}{\textsc{stands}}
\newcommand{\vR}{\textsc{stands w/ repair}}
\newcommand{\vF}{\textsc{fails}}
{\small\begin{tabular}{@{}llp{7.6cm}@{}}
\hline
Item & Verdict & Precise point\\ \hline
Lemma 1.1 & \vS & $\Tr(t_\delta\delta)=4g_Q$, $V(t_\delta)=2g_Q$; verified numerically\\
Remark 1.2 & \vR & right conclusion, inverted parenthetical: $\SLD(t)=\frac12\delta$ (R1); the slip is S10~(C4)'s (R2)\\
Lemma 2.1 & \vS & both directions are the same two inequalities\\
Lemma 2.2 & \vS & (1)--(3) correct for the polar $V$; abstract overclaims (R3)\\
Remark 2.3 & \vS & shift counterexample reproduced\\
Lemma 2.4, Prop.~2.5 & \vS & recomputed blockwise, exact at finite $s$\\
Def.~3.1 & \vR & no semigroup (R6); $N_1$ form vs operator (R13)\\
Lemma 3.2 & \vS & exact feasibility at finite $s$ verified from scratch\\
Lemma 3.3 & \vS & imported from S10 Lemma 11.1, correctly cited\\
Prop.~3.4 & \vS & all blocks, all signs recomputed; form statement only (R4)\\
Lemma 3.6 & \vR & cone is not a linear space and not closed\\
Thm.~3.9 (upper bd) & \vS & given H1; constructive, no circularity\\
Cor.~3.10 & \vS & class of $\mathcal G$ checked against S10 Thm.~6.2\\
Thm.~3.11(1) & \vR & Lumer--Phillips needs a semigroup (R6)\\
Thm.~3.11(2) & \vR & algebra correct, vacuous on $C_c^\infty$ (R5)\\
\textbf{Thm.~3.11(3)} & \vF & $+\lambda D_\chi|_D$ is anti-symmetric but not admissible; the claimed cone equality gives $T=0$ (R7)\\
Thm.~3.11(4) & \vS & bounded $\Rightarrow\sigma=0$; independent of (1)'s proof\\
Rem.~3.12, 3.13 & \vS & boundary form and over-compression recomputed\\
Prop.~3.14 & \vS & $w_{ik}\ge1$; block bounds\\
Rem.~3.15 & \vS & honest and correct about the essential-spectrum obstruction\\
H2, Thm.~3.17 & \vS & correct given H2; but H2 is the lower bound in disguise (R8)\\
Rem.~3.18 & \vR & needs $\Delta_z(t)\ge0$ (R9)\\
Def.~4.1 & \vR & $M_*,\tau_*$ are forms, not operators (R11)\\
Lemma 4.2 (formula) & \vS & recomputed; verified numerically to $10^{-9}$\\
\textbf{Lemma 4.2 (last clause)} & \vF & derivative $=2\lambda(1-\theta)g_Q(\dot\delta)>0$ along $\mathcal G=-\lambda D_\chi$ (R10)\\
Prop.~4.3 & \vR & correct on the linear class; test direction must be admissible (R10)\\
Thm.~4.4 (a)$\Leftrightarrow$(b),(c) & \vS & both directions airtight; no counterexample in $60\,000$ random trials\\
Thm.~4.4 (c)$\Leftrightarrow$(d) & \vR & incomplete: the one-sided boundary directions are not covered (R7, R10)\\
Cor.~4.5 & \vR & gain may vanish in the continuum limit (R14)\\
Prop.~4.6, Rem.~4.7 & \vS & limits and the conjugation by $X_c$ verified verbatim\\
\S5 (recipe for F2) & \vS & correct; add the rigidity arbiter and the cap scan (R14)\\
Lemma 6.1, Prop.~6.2 & \vS & $y$, $w=-4\sinh^2\frac y2$, $K$ all recomputed\\
Rem.~6.3 & \vS & numbers superseded (R12)\\
\textbf{Rem.~6.4, closed-form box} & \vF & diagnosis and window refuted by F3 (R12)\\
\S7 item 1 & \vR & ``co-isometry iff $1\notin\sigma_p(N)$'' (R3)\\
\textbf{\S7 item 4} & \vF & same as Thm.~3.11(3) (R7)\\
\hline
\end{tabular}}

\section{Required repairs}\label{sec:repairs}

\begin{enumerate}[leftmargin=2.6em,itemsep=1pt,label=\textbf{R\arabic*.}]
\item Remark~1.2: the parenthetical must read $\SLD(t)=\frac12\delta$ (equivalently
 $2\SLD(t)=\delta$), not $\SLD(t)=2\delta$.
\item For the S10 errata log (not this document): S10~(C4) is internally inconsistent ---
 with its own eq.~(5) $\SLD(t)=\delta$ the correct relations are $\Tr(t_\delta\delta)=4g_Q$,
 $V(t_\delta)=2g_Q$.  All S10 conclusions survive (the Legendre ratio and every sign
 condition are homogeneous).  F2 must state which normalisation its ``units of
 $g_Q(\dot\delta)$'' use.
\item Lemma~2.2(1), the abstract and \S7 item~1: ``$V$ co-isometry iff $1\notin\sigma_p(N)$''
 holds for the \emph{polar} $V$; the existence of \emph{some} co-isometric factor is
 $\dim\ker(1-N)\le\dim\ker X$.
\item Prop.~3.4: move ``as sesquilinear forms on $C_c^\infty$'' into the statement and add
 that H1 is what upgrades it to $g_Q$.
\item Define $\sigma$ by the boundary form $\langle\sigma g,g\rangle:=\Re\langle\zeta^*g,g\rangle$
 on $\operatorname{dom}\zeta^*$; ``$\frac12(\zeta+\zeta^*)$'' is not defined there.
\item Def.~3.1(1): require $(U_s)$ to be a $C_0$-semigroup of isometries (or assume the
 adjoint family weakly differentiable); without it Thm.~3.11(1) has no proof.
\item \textbf{Thm.~3.11(3) and \S7 item~4:} replace the cone \emph{equality} by the
 inclusion $\mathcal C\subseteq\{\Aop(\mathcal G,N_1',Y_1)\}$ with $\mathcal G$
 anti-symmetric \emph{and} an isometry generator (equivalently: $\mathcal G^*$ dissipative),
 $N_1'\ge0$ a form; delete ``(indeed unitary) $e^{s\mathcal G}$''.  State that the
 admissible $\mathcal G$ form a one-sided cone and that the enlarged cone would give $T=0$.
\item Say explicitly that H2 is not a regularity assumption but the lower bound in disguise
 (it assumes existence, the scaled form, boundedness and uniformity).
\item Remark~3.18: the weak-duality bound needs $\Delta_z(t)\ge0$.
\item \textbf{Lemma~4.2 (last clause), Prop.~4.3, Thm.~4.4(d):} restrict ``the derivative
 vanishes'' to the linear class $\overline{\mathcal Y}$, and add the missing one-sided
 conditions $-b(\delta_*,[Q,\zeta])\ge0$ for the boundary-moving isometry generators.
 Record that they hold for the over-compression direction ($=2\lambda(1-\theta)g_Q(\dot\delta)>0$),
 so the completed criterion is (a) \emph{plus} an endpoint condition.
\item Define $M_*,\tau_*$ as quadratic forms through $b(\delta_*,\cdot)$ (eq.~\eqref{eq:formdef}),
 state the domain $\mathfrak D$, and read (a) as a form inequality; $\lambda_{\min}$ is then
 an infimum that may be $-\infty$.
\item \S6.3: replace $0.30\lesssim\theta\lesssim0.40$ by $\theta=0.384\pm0.003$ (F3),
 replace the ``$y\gg1$ / $h$ cannot resolve $e^{-y}$'' diagnosis (the frames agree at
 matched $h$; the $h$-dependence is bulk; S10's $0.300$ is a taper-suppressed lower bound),
 and record that $3/\pi^2$, $1/3$ and $1-2/\pi$ are now excluded.  $1-\theta=0.616\pm0.003$.
\item Def.~3.1(2) vs Def.~3.2: $N_1$ must be a bounded operator for the scaled family to
 exist; ``quadratic forms'' is inconsistent with $(1-sN_1)^{1/2}$.
\item Cor.~4.5 and \S5: strictness $\kappa_{\rm opt}<1-\theta$ needs one admissible $\psi$
 with $\Delta<0$ \emph{and} $g_Q(d)<\infty$; ask F2 for $\sup_\psi\Delta^2/g_Q(d)$ as a
 function of $L$ and of the weight cap, together with
 $\|\mathrm{AntiHerm}((t_*Q)_{DD})\|$ as the M\"obius-rigidity arbiter.
\item \S4 preamble: define $\delta_*$ as the $g_Q$-orthogonal residual of $-\dot\delta$ on
 $\overline{\mathcal Y}$ (always exists) instead of positing a minimiser $\mathcal G_*$.
\end{enumerate}

\section{Summary}\label{sec:summary}
\begin{enumerate}[leftmargin=1.8em,itemsep=0pt]
\item The exact chart of $\cF$ (Lemmas~2.1--2.4, Prop.~2.5) and the exact feasibility of
 the scaled family at finite $s$ (Lemma~3.2) are correct and are the document's strongest
 contribution: every competitor is a genuine channel, so the upper bound is constructive.
\item Prop.~3.4 was recomputed block by block: $\Aop_{AD}=-Q_{AD}\zeta+\frac12Q_{AD}N_1$,
 $\Aop_{DD}=-(\zeta^*Q_{DD}+Q_{DD}\zeta)+\frac12\{N_1,Q_{DD}\}-Y_1$, all signs, the
 hermitisation and the factor $\frac12$ confirmed.  It is a \emph{form} identity.
\item Thm.~3.11(3) is false: $+\lambda D_\chi|_D$ is anti-symmetric but generates an
 outward flow, is not an isometry generator, and the claimed cone equality would force
 $T=0$.  Only the inclusion survives; the admissible anti-symmetric generators are a
 one-sided cone.  Item (4) (bounded $\Rightarrow$ nothing new) stands.
\item Consequently the last clause of Lemma~4.2 is false ($2\lambda(1-\theta)g_Q(\dot\delta)>0$
 along the over-compression direction) and Thm.~4.4(d) is incomplete: the criterion
 $M_*\ge0$, $M_*\ge\tau_*$ must be supplemented by an endpoint condition.
\item The derivative formula $\frac12[\Tr(N_1M_*)-\Tr(Y_1\tau_*)]$ and the equivalence
 (a)$\Leftrightarrow$(c) are correct; $60\,000$ random non-commuting $2\times2$/$3\times3$
 trials produced no counterexample, as the proof predicts.
\item $M_*$ and $\tau_*$ are in general only quadratic forms ($t_*\notin\Stwo$), so
 ``$M_*\ge0$'' needs a domain, $\lambda_{\min}$ may be $-\infty$, and Cor.~4.5's strict
 improvement needs a violating direction of \emph{finite} $g_Q$-norm.  F2's violations
 (gain $\sim10^{-10}g_Q(\dot\delta)$, collapsing by eight orders when the weight cap is
 raised, while the violation itself grows) therefore do not establish
 $\kappa_{\rm opt}<1-\theta$; F2's own ``artefact-suspect'' verdict on T2 is the right one.
\item Prop.~4.6 and Remark~4.7 are correct: S10's $\Sigma_2\ge0$ is the $N_1$ condition,
 $\Sigma_1\ge0$ the $Y_1$ condition, $\Sigma_1=\Sigma_1^*$ the rotation condition.
\item Prop.~6.2 (modular frame) is correct in all three ingredients; the $\theta$-diagnosis
 of Remark~6.4 is superseded by F3 ($\theta=0.384\pm0.003$, $1-\theta=0.616\pm0.003$).
\item H1 and H2 are honestly declared and the lower-bound discussion is honest about the
 essential-spectrum obstruction; six undeclared assumptions are listed in \S\ref{sec:hyp}.
 No circularity with S10 or E1.
\end{enumerate}

\section{Summary}\label{sec:summary}
\end{document}
